\chapter{Introduction}

Large language models can answer open-ended questions fluently even when
the answer they produce is unsupported or simply wrong. Such fluent
mistakes are commonly called \emph{hallucinations}, and in the
following we use this term as an umbrella for several distinct failure
modes (misinformation, confabulation, systematic reasoning errors).
A more careful taxonomy is given in
Chapter~\ref{ch:background}; for now the term refers to any answer the
model produces that is not supported by the input or by the facts.
The practical problem this raises is not merely that an answer is wrong,
but that a useful system should also be able to indicate \emph{when} its
answer is unreliable so that a downstream user, or an automated
verification step, can act on it.

Empirical evidence confirms the scale of the risk. Under adversarial
prompting, hallucination rates in clinical decision-support settings have
been measured as high as
50--82\,\%~\cite{omar_adversarial_hallucination_clinical_2025}; this is an
adversarial upper bound rather than a typical deployment rate, but
factuality failures remain common enough in ordinary use that the broader
landscape of deployed
LLMs~\cite{augenstein_factuality_challenges_2024} presents substantial
challenges for automated verification. The precise rate matters less here
than the structural problem it creates.
The central technical difficulty is that free-form natural language has
many valid surface forms. Two
sampled answers may use very different words while expressing the same
underlying claim, and two similar-looking answers may make incompatible
claims.

This thesis is closely based on the work of Farquhar et
al.~\cite{semantic_entropy_nature_2024}, who introduced \emph{semantic
entropy} as a training-free signal for hallucination risk. The basic
idea is taken directly from that paper: instead of measuring how much
the strings produced by the model vary, one tries to measure how much
the underlying \emph{meanings} of those strings vary. If repeated samples
from the model mostly express the same meaning, the model may be
uncertain only about phrasing; if the samples split across genuinely
different meanings, the model is semantically uncertain and the answer
should be treated as riskier, following the semantic uncertainty
framing of Kuhn et al.~\cite{kuhn_semantic_uncertainty_2023} and
Farquhar et al.~\cite{semantic_entropy_nature_2024}. The
goal of the present thesis is not to invent this idea, but to
re-implement it under a smaller, locally reproducible setup and to study
how reliably it works in that setting.

\section{Research Aim}

The main research question is whether semantic entropy provides a
better signal for hallucination risk and selective abstention than
surface-form uncertainty, and how that comparison changes when the
output shifts from short answers to paragraph-length generation. The
study covers TriviaQA, NQ-Open, and SVAMP as short-answer benchmarks
and a FactScore-style biography task as the long-form regime ---
evaluated both whole-paragraph and at the level of individual claims
--- across three open generators and three decoding seeds. It therefore
targets \emph{confabulation}-like failures, where the model's answer
varies semantically across samples drawn under the same prompt, while
also testing the harder case where surface variation is nearly
universal. It does not claim that semantic entropy detects every type
of hallucination, and in particular it cannot flag high-confidence
systematic errors where the model is consistently wrong, a failure
mode discussed by
Kalai et al.~\cite{kalai_why_language_models_hallucinate_2025} and
Simhi et al.~\cite{simhi_choke_2025}.

The working hypothesis from Farquhar et
al.~\cite{semantic_entropy_nature_2024} is that discrete semantic
entropy computed over bidirectional
entailment-based clusters achieves a higher AUROC --- a ranking metric
for placing incorrect answers above correct ones --- than surface-form
entropy computed over sampled answer strings on the same prompts. A
secondary hypothesis is that the improvement, if present, is at least
directionally consistent across the short-answer benchmarks, despite
their very different answer formats (open-domain factual recall vs.\
short numeric arithmetic). The biography task tests the complementary hypothesis
that meaning-aware uncertainty methods separate more clearly once raw
surface diversity ceases to be informative.

\section{Study Structure}

The study separates cheap surface-form baselines from the semantic
clustering condition that most directly follows Farquhar et
al.~\cite{semantic_entropy_nature_2024}.

\textbf{Surface-form baselines.} For each prompt, the
generation model is sampled several times, the answers are lightly
normalized, and two answers are grouped together only if their
normalized strings are identical. We call this \emph{exact-match}
clustering. A deliberately naive companion score counts raw strings
before normalization. These baselines make the full pipeline
auditable end to end. It does \emph{not} implement true semantic
entropy in Farquhar et al.'s sense~\cite{semantic_entropy_nature_2024}, because
exact-match grouping ignores paraphrases; throughout the thesis the
exact-match number is therefore reported as a surface-form baseline rather
than as semantic entropy.

\textbf{Replication condition.} The main semantic-entropy condition repeats
the same model--dataset grid on the NHR@FAU
high-performance computing cluster and replaces exact-match clustering
with a bidirectional Natural Language Inference (NLI) clustering step
using a DeBERTa-v3-base cross-encoder. This is the configuration that
corresponds most directly to the method of Farquhar et al.

\textbf{Additional analyses.} Separate sections report judge sensitivity
checks with DeBERTa-v3-large and Qwen2.5-72B-Instruct, a continuous
kernel-based variant called Kernel Language
Entropy~\cite{nikitin_kernel_language_entropy_2024}, and held-out
Semantic Entropy Probes. These analyses extend the replication question;
they are not part of the baseline.

\section{Contributions}

The substantive method comes from Farquhar et
al.~\cite{semantic_entropy_nature_2024}; the contribution of this thesis
is to re-implement it under a smaller, fully local setup and to use that
setup to ask a question the original aggregate evaluation does not
isolate --- how the value of meaning-level clustering scales as outputs
lengthen. This work therefore sets out to deliver, mainly through
implementation, evaluation, and scope choices:

\begin{enumerate}
  \item an implementation of discrete semantic entropy for short-answer
    and long-form hallucination detection, covering answer normalization,
    exact-match clustering, bidirectional NLI clustering with DeBERTa-v3,
    entropy scoring, and correctness labeling;
  \item an experimental design that separates surface-form baselines,
    the NLI-clustered replication condition, and additional analyses,
    with Slurm templates and reproducibility logging for each run;
  \item an auditable JSONL output schema recording sampled answers,
    semantic clusters, cluster representatives, entropy scores, and
    correctness labels for every evaluated prompt;
  \item an evaluation framework for selective prediction, using AUROC,
    rejection-accuracy curves, and their area summary (AURAC),
    applicable to any uncertainty estimator that writes a score field to
    the shared schema;
  \item extension baselines and sensitivity checks covering Semantic
    Entropy Probes (hidden-state classifiers), Kernel Language Entropy,
    P(True), in- and out-of-distribution embedding regression, and
    semantic-judge comparisons across DeBERTa-v3-base, DeBERTa-v3-large,
    DeBERTa-v2-xlarge-MNLI, and Qwen2.5-72B-Instruct;
  \item a controlled short-versus-long-form comparison that places the
    FactScore-style biography task alongside the short-answer
    benchmarks --- including a claim-level long-form extension --- and
    exposes the regime where meaning-aware methods separate cleanly from
    surface-form baselines.
\end{enumerate}

\section{Thesis Structure}

Chapter~\ref{ch:background} reviews the different forms of
LLM hallucination, the role of uncertainty estimation, the notion of
semantic equivalence, and related semantic uncertainty methods.
Chapter~\ref{ch:semantic_entropy} defines the discrete semantic entropy
method on which the empirical work is based.
Chapter~\ref{ch:experimental_design} turns the method into a reproducible
experimental design, and Chapter~\ref{ch:hpc} documents the
implementation and HPC execution. Chapter~\ref{ch:extensions}
describes Kernel Language Entropy as a completed posthoc check and
Semantic Entropy Probes as a held-out probe extension.
Chapter~\ref{ch:results} reports the empirical results,
Chapter~\ref{ch:discussion} discusses interpretation and limitations,
and Chapter~\ref{ch:conclusion} states the final conclusion.
